\cleardoublepage

\chapter*{Abstract\markboth{Abstract}{Abstract}}

Radio coverage planning for tactical 5G networks ---deployed in emergency scenarios,
military operations or temporary events--- requires simulation tools capable of accurately
modelling electromagnetic propagation over real terrain. Existing commercial solutions
(Xirio Online, Atoll) impose annual licence costs exceeding \texteuro2,000, restricting
access for budget-constrained organisations and academic environments.

This Final Degree Project designs and implements a proof-of-concept web-based digital twin
for coverage visualisation of 5G networks operating in the FR1 band (sub-6\,GHz). The
system integrates three components in a three-tier architecture: (i)~a MATLAB simulation
engine applying the Longley-Rice Irregular Terrain Model (ITM) propagation model over
30\,m resolution SRTM topographic data; (ii)~a REST backend built with Spring Boot~3.2
(Java~17) that serves simulation results through an OpenAPI-documented API; and (iii)~a
single-page web application in Angular~21 that renders interactive coverage heatmaps on
Leaflet-based cartography.

Three reference scenarios have been defined ---suburban (URJC campus, 3.5\,GHz), rural
mountainous (Sierra de Guadarrama, 700\,MHz) and dense urban (central Madrid,
3.5\,GHz)--- to validate model coherence against free-space theoretical bounds. The
prototype includes visual simulation comparison features, dynamic RSRP filtering and base
map selection. The backend incorporates 13 unit tests, path traversal security validation
and GZIP compression.

The outcome is an accessible, extensible and open-source tool that enables
telecommunications engineers to perform preliminary 5G coverage assessments without
reliance on commercial licences.

\vspace{1em}
\textbf{Keywords:} digital twin, 5G NR, radio coverage, Longley-Rice, SRTM, Spring Boot,
Angular, Leaflet, RSRP, tactical networks.
